\section{Explicit electrovacuum seeds}
\label{sec:construction}

We first solve the momentum and Maxwell divergence constraints explicitly.
The electromagnetic fields begin at a large radius; one scalar equation will
then determine the metric.  Throughout the construction,
\[
 \Omega=\{x\in\R^3:r=|x|\ge1\},\qquad S=\{r=1\},
\]
with the Euclidean metric $\delta$ and standard orientation.  Write
$(r,\theta,\phi)$ for spherical coordinates, with $\theta$ the colatitude and
$\eta=\partial_\phi$ the generator of the standard rotation action.
Fix a smooth function $\chi:\R\to[0,1]$ that is zero on $(-\infty,1]$ and
one on $[2,\infty)$.  For example, one may take
\[
 \chi(t)=\frac{\zeta(t-1)}{\zeta(t-1)+\zeta(2-t)},\qquad
 \zeta(t)=\begin{cases}e^{-1/t},&t>0,\\0,&t\le0.\end{cases}
\]
For $s\in[0,1]$ and $L>4$, put $c(r)=\chi(r/L)$.  Constants $C,C_j$ below
depend only on this fixed cutoff and the indicated derivative order, and
are independent of $s$ and $L$ above a fixed lower threshold.

Define the scalar potentials, vector fields, and symmetric covariant tensor
by
\begin{equation}
\label{eq:seeds}
\begin{aligned}
 f&=sc(r)\sin^2\theta,& h&=sc(r)\sin^2\theta\cos\theta,\\
 \iota_{\widetilde E}\dd V_\delta&=\dd(f\,\dd\phi),&
 \iota_{\widetilde B}\dd V_\delta&=\dd(h\,\dd\phi),\\
 \sigma&=-\frac{s^2c(r)^2\sin^4\theta}{r^2}
       (\dd r\otimes\dd\phi+\dd\phi\otimes\dd r).
\end{aligned}
\end{equation}
Here contraction with the volume form uniquely determines each vector
field.  The coordinate formulas in \eqref{eq:seeds} extend smoothly across
the rotation axis, as \Cref{lem:seeds} verifies.

\begin{lemma}[Seed identities and bounds]
\label{lem:seeds}
The fields $\widetilde E,\widetilde B$ and tensor $\sigma$ are smooth on
$\Omega$, invariant under the standard effective $U(1)$ action, and vanish
on $\{1\le r\le L\}$.  They satisfy
\begin{equation}
\label{eq:seed-constraints}
 \operatorname{div}_\delta\widetilde E
 =\operatorname{div}_\delta\widetilde B=0,\qquad
 \tr_\delta\sigma=0,\qquad
 \operatorname{div}_\delta\sigma
 =2\dd V_\delta(\widetilde E,\widetilde B,\cdot).
\end{equation}
For every Cartesian multi-index $\alpha$ of order $j\ge0$,
\begin{equation}
\label{eq:seed-bounds}
\begin{aligned}
 |\partial^\alpha\widetilde E|_\delta
 +|\partial^\alpha\widetilde B|_\delta
 &\le C_j s r^{-2-j}\mathbf1_{\{r\ge L\}},\\
 |\partial^\alpha\sigma|_\delta
 &\le C_j s^2 r^{-3-j}\mathbf1_{\{r\ge L\}}.
\end{aligned}
\end{equation}
Moreover, $\operatorname{div}_\delta\sigma$ is supported in
$\{L\le r\le2L\}$, with
\[
 |\partial^\alpha(\operatorname{div}_\delta\sigma)|_\delta
 \le C_j s^2r^{-4-j}\mathbf1_{\{L\le r\le2L\}}.
\]
\end{lemma}

\begin{proof}
Let $X=(x^1,x^2,x^3)$, $T=(-x^2,x^1,0)$, and
$w=(x^1)^2+(x^2)^2$, with $\flat$ denoting Euclidean metric dual in the
following formulas.  The primitives and tensor have the Cartesian expressions
\begin{equation}
\label{eq:cartesian-seeds}
\begin{aligned}
 f\,\dd\phi&=\frac{sc}{r^2}(x^1\dd x^2-x^2\dd x^1),\\
 h\,\dd\phi&=\frac{scx^3}{r^3}(x^1\dd x^2-x^2\dd x^1),\\
 \sigma&=-\frac{s^2c^2w}{r^7}
   (X^\flat\otimes T^\flat+T^\flat\otimes X^\flat).
\end{aligned}
\end{equation}
These are smooth on $r\ge1$, with no denominator involving distance to
the axis.  They also show rotation invariance and, since $X\cdot T=0$,
trace-freeness.  The exact flux forms are closed, so
$\dd(\iota_{\widetilde E}\dd V_\delta)=
(\operatorname{div}_\delta\widetilde E)\dd V_\delta=0$, and likewise for
$\widetilde B$.

For the momentum identity we compute off the axis and then use smoothness.
Since $\dd V_\delta=r^2\sin\theta\,\dd r\wedge\dd\theta\wedge\dd\phi$,
the spherical \emph{coordinate} components are
\begin{equation}
\label{eq:field-components}
\begin{aligned}
 \widetilde E^r&=\frac{2sc\cos\theta}{r^2},&
 \widetilde E^\theta&=-\frac{sc'\sin\theta}{r^2},&
 \widetilde E^\phi&=0,\\
 \widetilde B^r&=\frac{sc(3\cos^2\theta-1)}{r^2},&
 \widetilde B^\theta&=-\frac{sc'\sin\theta\cos\theta}{r^2},&
 \widetilde B^\phi&=0.
\end{aligned}
\end{equation}
In particular an orthonormal polar component is $r$ times the displayed
polar component.  The covector $\dd V_\delta(\widetilde E,\widetilde B,\cdot)$
has only an azimuthal component, namely
\begin{equation}
\label{eq:seed-cross-product}
 \dd V_\delta(\widetilde E,\widetilde B,\partial_\phi)
 =\frac{f_rh_\theta-f_\theta h_r}{r^2\sin\theta}
 =-\frac{s^2cc'\sin^4\theta}{r^2}.
\end{equation}
To check all three tensor-divergence components, put
$a=\sigma_{r\phi}=-s^2c^2\sin^4\theta/r^2$ and use the coordinate identity
\[
 (\operatorname{div}_\delta\sigma)_i
 =\frac1{\sqrt{\det\delta}}\partial_j
    \bigl(\sqrt{\det\delta}\,\sigma_i{}^j\bigr)
   -\frac12\sigma^{jk}\partial_i\delta_{jk}.
\]
The last term vanishes: the metric is diagonal in these coordinates,
whereas $\sigma^{jk}$ has only $r\phi$ and $\phi r$ entries.  Thus
\begin{equation}
\label{eq:all-divergence-components}
\begin{aligned}
 (\operatorname{div}_\delta\sigma)_r
  &=\frac1{r^2\sin\theta}\partial_\phi
       \left(\frac a{\sin\theta}\right)=0,\\
 (\operatorname{div}_\delta\sigma)_\theta&=0,\\
 (\operatorname{div}_\delta\sigma)_\phi
  &=\frac1{r^2\sin\theta}\partial_r(r^2\sin\theta\,a)
    =a_r+\frac{2a}{r}
    =-\frac{2s^2cc'\sin^4\theta}{r^2}.
\end{aligned}
\end{equation}
This proves \eqref{eq:seed-constraints}.  The last expression also proves
the stated support property; for $r\ge2L$ both fields are radial and
parallel.

Finally, for $j\ge1$ the Cartesian derivatives of $c$ obey
\[
 |\partial^\alpha c|\le C_j r^{-j}
   \mathbf1_{\{L\le r\le2L\}}.
\]
Indeed, all such derivatives are supported where $r$ is comparable to
$L$, and differentiating $\chi(r/L)$ gives this bound.  The coefficients
of the two uncut primitive one-forms in \eqref{eq:cartesian-seeds} are
smooth homogeneous functions of degree $-1$; exterior differentiation
therefore gives the field bounds in \eqref{eq:seed-bounds}.  The uncut
Cartesian tensor coefficients are smooth homogeneous functions of degree
$-3$, which gives the tensor bounds by the same product rule.  Smooth
cutoff jets give the estimates also at $r=L$ and $r=2L$.
\end{proof}

\subsection{Reduction to a scalar equation}

For a positive function $u$, define the physical data by
\begin{equation}
\label{eq:conformal-data}
 g=u^4\delta,\qquad K=u^{-2}\sigma,\qquad
 \mathcal E=u^{-6}\widetilde E,\qquad
 \mathcal B=u^{-6}\widetilde B.
\end{equation}
The next calculation fixes the conformal weights and the normalization
of the scalar equation directly.

\begin{lemma}[Conformal electrovacuum equations]
\label{lem:conformal}
For every smooth positive $u$, the data \eqref{eq:conformal-data} satisfy
\[
 \tau=\tr_g K=0,\qquad
 \operatorname{div}_g\mathcal E=\operatorname{div}_g\mathcal B=0,
 \qquad
 \operatorname{div}_gK=2(\mathcal E\times_g\mathcal B)^\flat_g.
\]
They satisfy the electrovacuum Hamiltonian constraint if and only if
\begin{equation}
\label{eq:scalar}
 -\Delta_\delta u=p u^{-7}+\ell u^{-3}\quad\hbox{on }\Omega,
\end{equation}
where, writing $t=\cos\theta$,
\begin{equation}
\label{eq:coefficients}
\begin{aligned}
 p&=\frac18|\sigma|_\delta^2
   =\frac{s^4c^4\sin^6\theta}{4r^6},\\
 \ell&=\frac14\bigl(|\widetilde E|_\delta^2
                  +|\widetilde B|_\delta^2\bigr)\\
   &=\frac{s^2}{4r^4}
      \left[c^2(9t^4-2t^2+1)+(rc')^2(1-t^4)\right].
\end{aligned}
\end{equation}
The nonnegative coefficients $p,\ell$ are smooth and rotation invariant.
For every Cartesian multi-index $\alpha$ of order $j\ge0$,
\begin{equation}
\label{eq:source-coefficient-bounds}
 |\partial^\alpha p|\le C_j s^4r^{-6-j}\mathbf1_{\{r\ge L\}},\qquad
 |\partial^\alpha\ell|\le C_j s^2r^{-4-j}\mathbf1_{\{r\ge L\}}.
\end{equation}
\end{lemma}

\begin{proof}
Since $\dd V_g=u^6\dd V_\delta$, the flux forms are unchanged:
\begin{equation}
\label{eq:conformal-fluxes}
 \iota_{\mathcal E}\dd V_g=\iota_{\widetilde E}\dd V_\delta,
 \qquad
 \iota_{\mathcal B}\dd V_g=\iota_{\widetilde B}\dd V_\delta.
\end{equation}
Their closedness proves the two Maxwell divergence constraints.  The
cross-product convention gives, for any vector $Z$,
\[
 (\mathcal E\times_g\mathcal B)^\flat_g(Z)
 =\dd V_g(\mathcal E,\mathcal B,Z)
 =u^{-6}\dd V_\delta(\widetilde E,\widetilde B,Z).
\]

Here is the tensor calculation in Cartesian coordinates.  Put $b=2\log u$,
so $g=e^{2b}\delta$ and $K=e^{-b}\sigma$.  The connection difference is
\[
 D^i_{jk}=\delta^i_j b_k+\delta^i_k b_j-\delta_{jk}b^i,
 \qquad b_j=\partial_jb,
\]
with indices on the right raised by $\delta$.  Symmetry and trace-freeness
give
\[
 \delta^{jk}D^l_{ki}\sigma_{lj}=0,\qquad
 \delta^{jk}D^l_{kj}=-b^l.
\]
Consequently the derivative of $e^{-b}$ cancels the second connection
contraction, yielding
\begin{equation}
\label{eq:conformal-divergence}
 \tr_gK=e^{-3b}\tr_\delta\sigma=0,\qquad
 (\operatorname{div}_gK)_i
 =e^{-3b}\delta^{jk}\partial_k\sigma_{ij}
 =u^{-6}(\operatorname{div}_\delta\sigma)_i.
\end{equation}
Together with \Cref{lem:seeds}, this is exactly the momentum constraint.

For completeness, substituting the same $D$ into
\[
 R_{ij}=\partial_kD^k_{ij}-\partial_jD^k_{ik}
          +D^k_{kl}D^l_{ij}-D^k_{jl}D^l_{ik}
\]
gives derivative terms $-b_{ij}-\delta_{ij}\Delta_\delta b$ and quadratic
terms $b_i b_j-\delta_{ij}|\nabla_\delta b|^2$.  Contracting therefore gives
\begin{equation}
\label{eq:conformal-curvature}
 R_g=e^{-2b}\bigl(-4\Delta_\delta b
                         -2|\nabla_\delta b|^2\bigr)
     =-8u^{-5}\Delta_\delta u.
\end{equation}
The norm identities are
\[
 |K|_g^2=u^{-12}|\sigma|_\delta^2,\qquad
 |\mathcal E|_g^2=u^{-8}|\widetilde E|_\delta^2,\qquad
 |\mathcal B|_g^2=u^{-8}|\widetilde B|_\delta^2.
\]
Thus the Hamiltonian constraint
$R_g-|K|_g^2=2(|\mathcal E|_g^2+|\mathcal B|_g^2)$ becomes
\eqref{eq:scalar}, with the coefficients in \eqref{eq:coefficients}.
Those explicit expressions follow from \eqref{eq:field-components} and
$|\sigma|_\delta^2=2a^2/(r^2\sin^2\theta)$.
Their smoothness, nonnegativity, and derivative bounds follow either from
their definitions as squared norms or directly from
\Cref{lem:seeds} and the product rule.
\end{proof}

We will solve \eqref{eq:scalar} with $u\to1$ at infinity and
$\partial_ru+u/2=0$ on $S$.  This boundary condition will make $S$ a future
MOTS.  A solution and the uniform estimates needed to control all other
surfaces are established next.
